%------------------------------------------------------------------------------------
%  RÉSUMÉ & ABSTRACT BILINGUE AVEC MOTS-CLÉS
%  Utilisation de \addchap pour une présence dans la table des matières sans numéro
%------------------------------------------------------------------------------------
\addchap{Résumé}

Ce document constitue à la fois un manuel de référence approfondi sur la typographie scientifique avec \LaTeX{} moderne et la vitrine exhaustive de la classe personnalisée \texttt{ClasseGuillaume}.
L'objectif central est de démontrer qu'il est possible d'allier une rigueur académique sans compromis à une esthétique visuelle équilibrée et aérée.
Toutes les facettes d'un mémoire ou d'un rapport d'ingénierie contemporain sont abordées : architecture modulaire des fichiers, compilation reproductible sous le moteur Lua\LaTeX{}, détection dynamique des polices de caractères, gestion fine des interlignes et élimination systématique des ruptures de lignes disgracieuses par microtypographie.

Le lecteur découvrira une mise en pratique détaillée des fonctionnalités exclusives de la classe :
le système collaboratif d'annotations marginales \texttt{\textbackslash todoH} avec troncature automatique et infobulles interactives, la gestion à la volée de géométries multiples (\texttt{\textbackslash geometrieWord}, \texttt{\textbackslash geometrieLarge}, \texttt{\textbackslash geometrieStandard}), la séparation stricte entre figures techniques et illustrations iconographiques dans des tables indépendantes, ainsi que l'harmonisation complète des tableaux sans barres verticales via le standard \textsf{booktabs}.
Ce modèle démontre qu'une publication de haut niveau repose sur une typographie maîtrisée où le gras, l'italique et les couleurs sont employés avec une stricte parcimonie.

\bigskip
\noindent\textbf{Mots-clés :} Lua\LaTeX{}, ClasseGuillaume, typographie scientifique, architecture modulaire, microtypographie, booktabs, siunitx, annotations collaboratives, scrlayer-scrpage.

\bigskip
\bigskip

\selectlanguage{english}
\addchap{Abstract}

This document serves as both an in-depth reference tutorial on modern scientific typesetting with \LaTeX{} and a comprehensive showcase of the custom \texttt{ClasseGuillaume} class.
Its core ambition is to demonstrate how uncompromising academic rigor can seamlessly pair with an elegant, balanced, and airy editorial design.
Every dimension of contemporary academic thesis and engineering report authoring is explored: modular source code architecture, reproducible compilation under the Lua\LaTeX{} engine, intelligent font detection, granular baseline stretch control, and microtypographical penalty optimization.

The reader is guided through practical implementations of the class's exclusive capabilities:
the collaborative marginal review system \texttt{\textbackslash todoH} featuring automatic message truncation and interactive PDF tooltips, on-the-fly page geometry toggling (\texttt{\textbackslash geometrieWord}, \texttt{\textbackslash geometrieLarge}, \texttt{\textbackslash geometrieStandard}), dedicated floating environments for illustrations alongside conventional figures, and publication-ready tables strictly adhering to the \textsf{booktabs} guidelines.
This template illustrates that high-end documents rely on disciplined typesetting where bold weights, italics, and chromatic accents are applied with deliberate parsimony.

\bigskip
\noindent\textbf{Keywords:} Lua\LaTeX{}, ClasseGuillaume, scientific typesetting, modular architecture, microtypography, booktabs, siunitx, collaborative annotations, scrlayer-scrpage.

\selectlanguage{french}
